\section{Theoretical Background}
\label{sec:background}

\subsection{Latent-variable models and the evidence lower bound}
A latent-variable model introduces unobserved variables $\z$ and defines
\begin{equation}
  p_\theta(\x) = \int p_\theta(\x \mid \z)\, p(\z)\, d\z .
\end{equation}
Maximum likelihood requires $\log p_\theta(\x)$, which is intractable for neural decoders. For any distribution
$q(\z \mid \x)$, Jensen's inequality gives the \emph{evidence lower bound}
\begin{equation}
  \log p_\theta(\x) \;\ge\; \E_{q(\z\mid\x)}\!\big[\log p_\theta(\x \mid \z)\big] - \KL{q(\z\mid\x)}{p(\z)}
  \;=\; \elbo(\x),
  \label{eq:elbo}
\end{equation}
with equality when $q(\z\mid\x)$ is the true posterior $p_\theta(\z\mid\x)$. Both models in this report maximise a
bound of this form. They differ in \emph{what} the latent variables are and \emph{whether} $q$ is learned.

\subsection{Variational Autoencoder}
\label{sec:bg-vae}

\paragraph{Generative model.} A latent code $\z\in\mathbb{R}^{d_z}$ is drawn from a standard normal prior and
decoded by a neural network into a diagonal Gaussian over pixels:
\begin{equation}
  p(\z) = \N(\z; \mathbf{0}, \I), \qquad
  p_\theta(\x\mid\z) = \N\!\big(\x;\, \mub_\theta(\z),\, \operatorname{diag}\sigb^2_\theta(\z)\big).
\end{equation}

\paragraph{Inference model (encoder).} The intractable posterior is approximated by an amortised Gaussian
\begin{equation}
  q_\phi(\z\mid\x) = \N\!\big(\z;\, \mub_\phi(\x),\, \operatorname{diag}\sigb^2_\phi(\x)\big),
\end{equation}
whose parameters are predicted by an encoder network in a single forward pass.

\paragraph{Reparameterisation.} To back-propagate through the sampling of $\z$, the sample is written as a
deterministic function of $\phi$ and an independent noise variable:
\begin{equation}
  \z = \mub_\phi(\x) + \sigb_\phi(\x)\odot\eps, \qquad \eps \sim \N(\mathbf{0},\I).
\end{equation}
The expectation in \cref{eq:elbo} is estimated with a single sample ($L=1$), as in the original paper.

\paragraph{The two ELBO terms in closed form.} With $D = 3\times32\times32 = 3072$ pixel dimensions and $d_z$
latent dimensions, the reconstruction term of one sample is
\begin{equation}
  -\log p_\theta(\x\mid\z) = \frac{1}{2}\sum_{i=1}^{D}\left[\log 2\pi + \log\sigma^2_{\theta,i}(\z)
  + \frac{\big(x_i-\mu_{\theta,i}(\z)\big)^2}{\sigma^2_{\theta,i}(\z)}\right],
  \label{eq:vae-recon}
\end{equation}
and the KL divergence between two diagonal Gaussians is analytic:
\begin{equation}
  \KL{q_\phi(\z\mid\x)}{\N(\mathbf{0},\I)} = \frac{1}{2}\sum_{j=1}^{d_z}\left(\mu_{\phi,j}^2 + \sigma_{\phi,j}^2
  - \log\sigma_{\phi,j}^2 - 1\right).
  \label{eq:vae-kl}
\end{equation}
Both sums run over \emph{all} dimensions before averaging over the batch, so the KL weight is exactly $\beta=1$. A
common implementation error averages \cref{eq:vae-recon} per pixel but sums \cref{eq:vae-kl}. That silently
multiplies the KL weight by $D=3072$ and causes posterior collapse; a unit test in the repository guards against it.

\paragraph{Generation.} A new image is obtained by drawing $\z\sim\N(\mathbf{0},\I)$ and outputting
$\mub_\theta(\z)$. This takes \emph{one} decoder evaluation.

\paragraph{Discrete data and bits/dim.} Pixels are integers $k\in\{0,\dots,255\}$. To make a continuous density a
valid bound on the discrete likelihood, we use uniform dequantisation:
\begin{equation}
  x = \frac{k+u}{128} - 1, \qquad u\sim\mathcal{U}[0,1),
\end{equation}
which maps every bin to an interval of width $w=2/256$ in $[-1,1)$. Since
$P(k) = \int_{\text{bin}} p(x)\,dx \ge w \exp\E_u[\log p(x)]$, the negative ELBO $\mathcal{B}$ (in nats) of an
image converts into a bound in bits per dimension:
\begin{equation}
  \mathrm{bpd} = \frac{\mathcal{B} + D\ln 128}{D\ln 2}.
  \label{eq:bpd}
\end{equation}

\subsection{Denoising Diffusion Probabilistic Model}
\label{sec:bg-ddpm}

\begin{figure}[t]
\centering
\begin{tikzpicture}[
  node distance=1.55cm,
  st/.style={circle,draw=lineC,thick,minimum size=0.95cm,inner sep=0pt,font=\small},
  fw/.style={-{Stealth[length=2.2mm]},thick,blue!60!black},
  bw/.style={-{Stealth[length=2.2mm]},thick,red!65!black},
]
\node[st,fill=midC] (x0) {$\x_0$};
\node[st,right=of x0] (x1) {$\x_1$};
\node[right=0.9cm of x1] (dots1) {$\cdots$};
\node[st,right=0.9cm of dots1] (xtm1) {$\x_{t-1}$};
\node[st,right=of xtm1] (xt) {$\x_t$};
\node[right=0.9cm of xt] (dots2) {$\cdots$};
\node[st,right=0.9cm of dots2,fill=latC] (xT) {$\x_T$};
\draw[fw] (x0) -- (x1);
\draw[fw] (xtm1) -- node[above,font=\footnotesize]{$q(\x_t\mid\x_{t-1})$} (xt);
\draw[bw] (xt) to[bend left=45] node[below,font=\footnotesize]{$p_\theta(\x_{t-1}\mid\x_t)$} (xtm1);
\draw[bw] (x1) to[bend left=45] (x0);
\node[below=0.15cm of x0,font=\scriptsize,align=center,yshift=-0.55cm] {data};
\node[below=0.15cm of xT,font=\scriptsize,align=center,yshift=-0.55cm] {$\approx\N(\mathbf{0},\I)$};
\node[above=0.35cm of x1,font=\scriptsize,blue!60!black,xshift=1.4cm] {fixed forward process (adds noise)};
\node[below=0.75cm of dots1,font=\scriptsize,red!65!black,xshift=1.6cm] {learned reverse process (removes noise)};
\end{tikzpicture}
\caption{The DDPM as a Markov chain. The forward process $q$ is fixed and gradually destroys the image; a single
network, shared across all $t$, parameterises the reverse process $p_\theta$, which generates images from noise.}
\label{fig:ddpm-chain}
\end{figure}

\paragraph{Forward (diffusion) process.} A fixed Markov chain adds Gaussian noise over $T=1000$ steps
(\cref{fig:ddpm-chain}):
\begin{equation}
  q(\x_t\mid\x_{t-1}) = \N\!\big(\x_t;\,\sqrt{1-\beta_t}\,\x_{t-1},\, \beta_t\I\big),
  \qquad \beta_t = 10^{-4} + \frac{t-1}{T-1}\,(0.02-10^{-4}).
\end{equation}
With $\alpha_t = 1-\beta_t$ and $\abar_t=\prod_{s=1}^{t}\alpha_s$, the marginal at any step is available in closed
form, so $\x_t$ can be sampled in one shot during training:
\begin{equation}
  q(\x_t\mid\x_0) = \N\!\big(\x_t;\,\sqrt{\abar_t}\,\x_0,\,(1-\abar_t)\I\big)
  \;\;\Longleftrightarrow\;\;
  \x_t = \sqrt{\abar_t}\,\x_0 + \sqrt{1-\abar_t}\,\eps,\quad \eps\sim\N(\mathbf{0},\I).
  \label{eq:q-marginal}
\end{equation}
With this schedule, $\abar_T \approx 4\times10^{-5}$, so $\x_T$ is almost exactly standard normal.

\paragraph{Forward-process posterior.} Conditioned on $\x_0$, the reverse of a forward step is a tractable Gaussian:
\begin{align}
  q(\x_{t-1}\mid\x_t,\x_0) &= \N\!\big(\x_{t-1};\,\mutil_t(\x_t,\x_0),\,\btil_t\I\big),\\
  \mutil_t(\x_t,\x_0) &= \frac{\sqrt{\abar_{t-1}}\,\beta_t}{1-\abar_t}\,\x_0
                        + \frac{\sqrt{\alpha_t}\,(1-\abar_{t-1})}{1-\abar_t}\,\x_t,
  \qquad
  \btil_t = \frac{1-\abar_{t-1}}{1-\abar_t}\,\beta_t .
  \label{eq:posterior}
\end{align}

\paragraph{Reverse (generative) process.} A neural network parameterises Gaussian transitions that start from pure
noise:
\begin{equation}
  p(\x_T)=\N(\mathbf{0},\I), \qquad
  p_\theta(\x_{t-1}\mid\x_t) = \N\!\big(\x_{t-1};\,\mub_\theta(\x_t,t),\,\sigma_t^2\I\big),
\end{equation}
with fixed variances $\sigma_t^2=\beta_t$ (``fixed-large'', used for CIFAR-10) or $\sigma_t^2=\btil_t$
(``fixed-small'').

\paragraph{Variational bound.} The negative ELBO of the chain decomposes into a sum of per-step terms:
\begin{equation}
\begin{split}
  -\log p_\theta(\x_0) \;\le\;\; &
  \underbrace{\KL{q(\x_T\mid\x_0)}{p(\x_T)}}_{L_T}
  \;+\; \sum_{t=2}^{T}\underbrace{\E_q\,\KL{q(\x_{t-1}\mid\x_t,\x_0)}{p_\theta(\x_{t-1}\mid\x_t)}}_{L_{t-1}}\\
  & \underbrace{-\,\E_q\log p_\theta(\x_0\mid\x_1)}_{L_0}.
\end{split}
\label{eq:vlb}
\end{equation}
$L_T$ has no parameters, and each $L_{t-1}$ is a KL divergence between two Gaussians. The final term $L_0$ uses a
discretised Gaussian decoder that integrates the density over each pixel's bin (width $2/255$, with the outermost
bins extended to $\pm\infty$), so the bound is a proper bound on the discrete likelihood.

\paragraph{Noise parameterisation and $L_{\text{simple}}$.} Substituting \cref{eq:q-marginal} into
\cref{eq:posterior} suggests predicting the noise instead of the mean:
\begin{equation}
  \mub_\theta(\x_t,t) = \frac{1}{\sqrt{\alpha_t}}\left(\x_t - \frac{\beta_t}{\sqrt{1-\abar_t}}\,\eps_\theta(\x_t,t)\right).
  \label{eq:eps-param}
\end{equation}
With this choice, $L_{t-1}$ becomes a weighted squared error on the noise. \citet{ho2020ddpm} drop the weights and
train with
\begin{equation}
  L_{\text{simple}}(\theta) = \E_{t\sim\mathcal{U}\{1..T\},\,\x_0,\,\eps}
  \Big[\big\|\eps - \eps_\theta\big(\sqrt{\abar_t}\,\x_0+\sqrt{1-\abar_t}\,\eps,\; t\big)\big\|^2\Big],
  \label{eq:lsimple}
\end{equation}
a reweighted ELBO that emphasises the harder, noisier steps and gives the best samples. In our implementation the
mean is computed through the predicted clean image
$\hat\x_0=(\x_t-\sqrt{1-\abar_t}\,\eps_\theta)/\sqrt{\abar_t}$, clipped to $[-1,1]$, followed by
$\mub_\theta=\mutil_t(\x_t,\hat\x_0)$. Without clipping this is algebraically identical to \cref{eq:eps-param}.

\paragraph{Ancestral sampling.} Generation runs the reverse chain from $\x_T\sim\N(\mathbf{0},\I)$:
$\x_{t-1} = \mub_\theta(\x_t,t) + \sigma_t\mathbf{z}$ with $\mathbf{z}\sim\N(\mathbf{0},\I)$ for $t>1$. This
takes $T=1000$ network evaluations per image.

\paragraph{DDIM sampling.} \citet{song2021ddim} show that the same trained network $\eps_\theta$ also defines a
deterministic, non-Markovian sampler on a subsequence $\tau_1<\dots<\tau_S$ of timesteps. For $\eta=0$:
\begin{equation}
  \x_{\tau_{i-1}} = \sqrt{\abar_{\tau_{i-1}}}\;\hat\x_0(\x_{\tau_i})
  + \sqrt{1-\abar_{\tau_{i-1}}}\;\eps_\theta(\x_{\tau_i},\tau_i).
  \label{eq:ddim}
\end{equation}
This trades quality for speed with \emph{no retraining}, which we use to measure the speed--quality frontier.

\subsection{A unifying view}
\label{sec:bg-unify}
\Cref{eq:vlb} is the ELBO of a \emph{hierarchical VAE} with latents $\x_1,\dots,\x_T$. Read this way, a DDPM is a
VAE with three special properties:
\begin{enumerate}
  \item the ``encoder'' $q(\x_{1:T}\mid\x_0)$ is fixed and has no parameters;
  \item every latent has the same dimension as the image, so nothing is compressed;
  \item the prior is matched by construction ($L_T\approx0$), so the KL to the prior costs nothing.
\end{enumerate}
The VAE makes the opposite choices: it has a single latent layer, a \emph{learned} encoder, a compressed latent
($d_z\ll D$), and a KL term the encoder must actively pay for. \Cref{sec:discussion} shows that these choices
explain every difference we measure.
